\section{相对运动的矢量关系}
\subsection{先分清三个速度的参照}
研究船、河水和河岸时，同一个“船速”可能指船对水的速度，也可能指船对岸的速度。最稳妥的记法是用下标写清参照：$\bm v_{A/B}$ 表示 $A$ 相对 $B$ 的速度，$\bm v_{B/O}$ 表示运动参考系 $B$ 相对地面 $O$ 的速度。对平动且坐标轴相互平行的参考系，位置满足 $\bm r_{A/O}=\bm r_{B/O}+\bm r_{A/B}$，同一时间下求导得
\begin{equation}\bm v_{A/O}=\bm v_{A/B}+\bm v_{B/O},\qquad \bm a_{A/O}=\bm a_{A/B}+\bm a_{B/O}.\label{eq:k2relative}\end{equation}
图\ref{k2theory-triangle}中的箭头首尾相接，表示式\eqref{eq:k2relative}的矢量和。先相加的是分量，不是速率。若求反向参照，应使用 $\bm v_{B/A}=-\bm v_{A/B}$。将两速度直接相减时，仍需先规定它们原来均以谁为参照。
\fig{k2theory-triangle}{平动参考系之间的速度三角形}{\begin{tikzpicture}
\draw[flow](0,0)--(3,0)node[midway,below]{$\bm v_{B/O}$};\draw[flow,blue](3,0)--(4.3,2)node[midway,right]{$\bm v_{A/B}$};\draw[flow,red](0,0)--(4.3,2)node[midway,above left]{$\bm v_{A/O}$};\end{tikzpicture}}

\subsection{运动合成与独立分解}
把一个二维运动分解到相互垂直的两轴，分运动共用同一个时间。分解后的初速度与加速度必须由原矢量投影得到，不能各自任意指定一个运动再拼合。在恒定重力场中，水平和竖直方向满足独立的标量方程，所以可先由竖直运动确定时间，再代入水平运动求位移。

从同一位置、同一时刻出发且加速度相同的两物体，位置之差满足 $\Delta\bm r=(\bm v_{20}-\bm v_{10})t$。共同加速度在相减时消失；若初始位置不同，还须保留初始相对位矢。这一结论能解释成束抛出的物体为何保持由初速度决定的几何形状，但不表示它们各自相对地面做匀速直线运动。

式\eqref{eq:k2relative}的加速度形式要求坐标轴不相对转动。对于旋转木马上的观察者，只能在某瞬间使用当地观察点的速度作速度差；若继续讨论有限时间的相对位移或加速度，须考虑参考系方向的转动。把某一瞬间的牵连速度当作全程常矢量，往往会错误地把圆周运动画成直线。


图\ref{k2deep:common}取同时从同一高度水平抛出的两物体，水平初速度分别为 $u_1,u_2$。同一时刻，它们下落的高度都为 $gt^2/2$，因此连线保持水平，长度为 $|u_2-u_1|t$。每个物体的地面轨迹仍是抛物线，而在随物体1平动、坐标轴不转动的系中，物体2沿水平直线匀速运动。
\fig{k2deep:common}{共同竖直加速度相消，同一时刻的相对位移保持水平}{\begin{tikzpicture}
\draw[axis](0,0)--(5.5,0)node[right]{$x$};\draw[axis](0,0)--(0,3.6)node[above]{$y$};
\draw[blue,thick,domain=0:2.2,samples=60]plot(\x,{3-.5*\x*\x});\draw[red,thick,domain=0:3.96,samples=70]plot(\x,{3-.5*(\x/1.8)*(\x/1.8)});
\fill(0,3)circle(2pt)node[left]{$A$};\fill[blue](1.8,1.38)circle(2pt)node[above left]{$P_1(t)$};\fill[red](3.24,1.38)circle(2pt)node[above right]{$P_2(t)$};
\draw[flow](1.8,1.38)--(3.24,1.38)node[midway,above]{$\Delta x$};\draw[flow](1.8,1.15)--(1.8,.5)node[left]{$\bm g$};\draw[flow](3.24,1.15)--(3.24,.5)node[right]{$\bm g$};
\node at(3.8,3.3){$u_2>u_1$};
\end{tikzpicture}}
图中的两点必须取同一时刻；若把两条轨迹上任意两个点相减，共同的 $gt^2/2$ 就不能抵消。若两物体落地时间不同，地面约束介入后它们不再具有相同加速度，相对匀速结论也随之终止。


\subsection{由相对方向反求风与雨}
无风竖直雨在行驶车辆中看来会向后倾斜，因为车辆速度从雨滴速度中被减去。取车辆向右为正，雨的地面速度为 $(u,-w)$，车速为 $V$，则雨对车速度为 $(u-V,-w)$。倾斜方向由 $u-V$ 的正负决定，倾角正切由水平、竖直分量之比给出。题目中的“风从东北吹来”描述来向，风速矢量却指向西南，作图前必须转换。

图\ref{k2theory-rain}把地面观察和车辆观察放在同一坐标方向下，蓝箭头多出的向后分量来自参照改变。若用雨痕反推相对速度，应明确将痕迹方向理想化为雨滴相对运动的投影；实际水滴在玻璃上的流动还受重力、黏附和气流影响。相对速度题所需的是这一运动学模型，不是对所有真实雨痕的无条件解释。
\fig{k2theory-rain}{车辆速度改变雨滴的相对方向}{\begin{tikzpicture}
\draw[axis](0,0)--(3.5,0)node[right]{$x$};\draw[flow](1.6,1.5)--(1.6,-.8)node[midway,right]{雨对地};\draw[flow,blue](1.6,1.5)--(-.5,-.8)node[left]{雨对车};\draw[flow](1.6,-.8)--(3,-.8)node[right]{$V$};\end{tikzpicture}}

\subsection{相对速度怎样进入穿越数量}
相对方向解决“从哪里来”，相对速度的法向分量解决“穿过多少”。设均匀分布的小粒子数密度为 $n$，平面面积为 $S$，其法向单位矢量为 $\bm n$。在时间 $\dd t$内，粒子相对平面的法向位移为 $(\bm v_{\text{粒}}-\bm v_{\text{面}})\cdot\bm n\,\dd t$，对应被扫过的薄柱体积为该位移绝对值乘 $S$。因此穿过一面平面的数量为
\[\dd N=nS\bigl|(\bm v_{\text{粒}}-\bm v_{\text{面}})\cdot\bm n\bigr|\dd t.\]
平行于平面的速度只改变交点位置，不改变这一面上的通量。若数密度或朝向变化，应沿时间积分；若只统计从某一侧入射，还须取指向该侧的正部分，不能把双向穿越全部计入。

对迎雨运动的竖直板，保持行驶距离时提高车速会增大单位时间内扫过的空气体积，却同时缩短运动时间。两个因素可能抵消，结果不一定随车速改变。对水平伞面，起作用的则是相对速度的竖直分量；板与伞的法向不同，不能只凭同一句“跑快些”作判断。

\fig{fv:k2:flux}{移动平面扫过的是相对法向薄柱，而非粒子的全部路程}{\begin{tikzpicture}[scale=.9]
\fill[blue!9](1,-1)rectangle(3,1.5);
\draw[thick](3,-1.5)--(3,2);\draw[dashed](1,-1.5)--(1,1.5);
\node at(3,2.5){随平面平移的坐标系};
\draw[->,blue,thick](0,.1)--(2,.9)node[midway,above left]{$\bm v_{\rm rel}$};
\draw[->](3,.1)--(4.5,.1)node[right]{$\bm n$};
\draw[<->](1,-.7)--(3,-.7)node[midway,below]{$|v_{{\rm rel},n}|\dd t$};\end{tikzpicture}}
图\ref{fv:k2:flux}先把观察区域随平面平移，粒子的速度相应减去平面速度。只有相对速度在法线 $\bm n$上的投影决定是否穿面，平行于面的分量只改变落点。故薄柱体积为 $S|v_{{\rm rel},n}|\dd t$。若两者法向速度相等，粒子可以在面前迅速横向移动，却不会穿过它；这能检查流量式中是否误用了总速率。


\section{匀变速直线运动的完整条件}
\subsection{三个常用关系怎样得到}
沿固定直线建立带正负的坐标，若加速度 $a$ 恒定，从 $t=0$ 的初值 $v_0,x_0$ 开始，逐次积分为
\begin{align*}
 v(t)-v_0&=\int_0^t a\dd\tau=[a\tau]_0^t=at,\\
 x(t)-x_0&=\int_0^t(v_0+a\tau)\dd\tau
 =\left[v_0\tau+\frac a2\tau^2\right]_0^t
 =v_0t+\frac12at^2.
\end{align*}
将第二式乘 $2a$，再用 $at=v-v_0$，得 $2a(x-x_0)=2v_0(v-v_0)+(v-v_0)^2=v^2-v_0^2$。这一步没有除以 $a$，故 $a=0$ 时仍有等式，但平方关系不能独自确定匀速方向和位移，应回到 $v=v_0$ 与 $x-x_0=v_0t$。最后一式没有时间信息，也不直接给速度正负，开平方后必须根据运动分支选根。

平均速度 $\overline v=(v_0+v)/2$ 来自速度图是直线；它不是一般变速运动的公式。若物体先运动后停住，随后保持静止，原来的 $v_0+at$ 只适用于停车之前。停车之后的位置保持不变，不能继续外推使汽车“倒着刹车”。求安全距离时，正是这种分段决定最大接近量发生在何时。

\subsection{分段过程与初值衔接}
若在时刻 $T$ 更换加速度，第二段应使用当地时间 $\tau=t-T$，并以第一段结束的位置和速度为新初值。譬如先以 $a_1$ 运动 $T$，再以 $a_2$ 运动 $T$，总位移为 $v_0T+a_1T^2/2+(v_0+a_1T)T+a_2T^2/2$。第二段继承的速度含有 $a_1T$，所以交换两段加速度的顺序通常会改变总位移。

图\ref{k2theory-piece}中，折点前后的速度相同而斜率不同，速度曲线下面积必须按每段自己的初速度计算。边界条件还包括限速、停车线、碰撞、到达斜面末端等事件；事件发生后，适用的方程可能改变。先列出有效时间区间，再解代数方程，能排除落在已经失效阶段的伪解。
\fig{k2theory-piece}{分段加速度改变斜率但保留速度的连续性}{\begin{tikzpicture}
\draw[axis](0,0)--(5.5,0)node[right]{$t$};\draw[axis](0,0)--(0,3.4)node[above]{$v$};\draw[world](0,.4)--(2,1.3)--(4,3);\draw[guide](2,1.3)--(2,0);\node[below]at(2,0){$T$};\node[below]at(4,0){$2T$};\node[above]at(1,1){$a_1$};\node[above left]at(3,2.2){$a_2$};\end{tikzpicture}}


图\ref{k2deep:order}比较先大后小与先小后大的两段加速度，每段都持续 $T$。两种过程的末速度都为 $v_0+(a_1+a_2)T$，但若 $a_1>a_2$，先大后小的速度曲线在中间阶段始终较高，因此总位移更大。两条图线之间的菱形面积为
\[
 \Delta x=(a_1-a_2)T^2.
\]
这个差值也可由两段位移式相减得到；图像直接说明了原因：较早获得的速度增量会在后续更长时间内贡献位移。
\fig{k2deep:order}{相同末速度不保证相同位移，加速顺序改变图下面积}{\begin{tikzpicture}
\draw[axis](0,0)--(5.6,0)node[right]{$t$};\draw[axis](0,0)--(0,3.8)node[above]{$v$};
\fill[blue!12](0,.5)--(2.4,2.5)--(4.8,3.3)--(2.4,1.3)--cycle;
\draw[blue,thick](0,.5)--(2.4,2.5)--(4.8,3.3);\draw[red,thick](0,.5)--(2.4,1.3)--(4.8,3.3);
\draw[guide](2.4,0)--(2.4,2.5) (4.8,0)--(4.8,3.3);\node[below]at(2.4,0){$T$};\node[below]at(4.8,0){$2T$};\node[left]at(0,.5){$v_0$};
\node[blue,above left]at(1.2,1.5){先 $a_1$};\node[red,below right]at(1.2,.9){先 $a_2$};\node[right]at(4.8,3.3){共同末速};
\end{tikzpicture}}


\subsection{等时间、等位移的不同平均}
匀加速直线运动中，第 $n$个相等时长 $\tau$对应区间 $[(n-1)\tau,n\tau]$。将两端代入同一位置函数再相减，得到
\begin{align*}
 \Delta x_n&=x(n\tau)-x((n-1)\tau)\\
 &=v_0[n-(n-1)]\tau+\frac a2[n^2-(n-1)^2]\tau^2\\
 &=v_0\tau+\frac a2(2n-1)\tau^2,\\
 \Delta x_{n+1}-\Delta x_n&=\frac a2[(2n+1)-(2n-1)]\tau^2=a\tau^2.
\end{align*}
由静止出发且全程不转向时，各段路程比为 $1:3:5:\cdots$；若含转向，带符号位移仍满足公式，路程比却须重新分段。这个条件不能省略。

同一匀加速阶段，时间中点的速度等于首末速度算术平均，而位移中点满足 $v_m^2=(v_0^2+v_1^2)/2$。前者由 $v(t)$线性，后者由 $v^2(x)$线性得到。把“中间时刻”误读成“走到中点”，会使两个看似相近的问题得到错误的平均值。分段更换加速度后，整段过程也不再自动具有这些单一阶段的中点性质。


图\ref{k2deep:midpoint}用两个不同横轴说明两个“中点”。左图的直线性属于 $v$ 对 $t$，所以时间中点直接平均两端速度；右图的直线性属于 $v^2$ 对 $x$，所以位移中点先平均速度平方，再按运动支路取平方根。图中取全程正向且 $v_1>v_0>0$，两者分别为
\[
 v_{\text{时中}}=\frac{v_0+v_1}{2},\qquad
 v_{\text{位中}}=\sqrt{\frac{v_0^2+v_1^2}{2}}.
\]
\fig{k2deep:midpoint}{在各自线性图上取中点，平均的是纵轴实际表示的量}{\begin{tikzpicture}[scale=.9]
\begin{scope}\draw[axis](0,0)--(4,0)node[right]{$t$};\draw[axis](0,0)--(0,3)node[above]{$v$};\draw[thick](0,.6)--(3.4,2.5);\draw[guide](1.7,0)--(1.7,1.55)--(0,1.55);\fill(1.7,1.55)circle(1.6pt);\node[below]at(1.7,0){$T/2$};\node[left]at(0,.6){$v_0$};\node[right]at(3.4,2.5){$v_1$};\end{scope}
\begin{scope}[xshift=5.5cm]\draw[axis](0,0)--(4,0)node[right]{$x$};\draw[axis](0,0)--(0,3)node[above]{$v^2$};\draw[thick](0,.4)--(3.4,2.5);\draw[guide](1.7,0)--(1.7,1.45)--(0,1.45);\fill(1.7,1.45)circle(1.6pt);\node[below]at(1.7,0){$(x_0+x_1)/2$};\node[left]at(0,.4){$v_0^2$};\node[right]at(3.4,2.5){$v_1^2$};\end{scope}
\end{tikzpicture}}
存在转向时，位移坐标可能重复，同一位置可对应正负两个速度；右图只能确定速度平方，不能替代支路判断。前后加速度不同时，两张图也不再整体为同一条直线，须按阶段分别使用。


\subsection{按加速度的自变量选择积分}
运动方程 $\dot v=a$没有改变，改变的是已知量的表达方式。若给 $a(t)$，先积分时间；若给 $a(x)$，利用 $\dot x=v$消去时间；若给 $a(v)$，在 $a(v)\ne0$的运动分支上写为 $\dd t=\dd v/a(v)$，进一步有 $\dd x=v\dd v/a(v)$。由此得到
\[t-t_0=\int_{v_0}^{v}\frac{\dd u}{a(u)},\qquad
x-x_0=\int_{v_0}^{v}\frac{u\dd u}{a(u)}.\]
这两式的积分上下限带有速度变化方向。当物体减速时，$\dd v$与 $a$可同为负，使时间仍增加；若将加速度大小直接代入分母，会把符号关系破坏。若某处 $a=0$，还需判断能否跨过该速度，不能在积分中任意穿过奇点。

同样一个恒减速度过程可以分别用时间、速度、位置积分求解，结果应一致。这种交叉检查能发现初值遗漏，也能表明哪一种变量使分段最少。求是否在有限距离停住，位置积分可能比求完整 $v(t)$更直接；求延迟反应的影响，时间分段则更自然。

\subsection{由位置相关的加速度求速度}
当已知 $a=a(x)$，而时间函数不易直接求出时，用链式法则 $a=\dd v/\dd t=(\dd v/\dd x)(\dd x/\dd t)=v\dd v/\dd x$。将两侧微分关系带初值积分，
\[
 \int_{v_0}^{v}u\dd u=\int_{x_0}^{x}a(\xi)\dd\xi,
 \qquad
 \left[\frac{u^2}{2}\right]_{v_0}^{v}=\frac{v^2-v_0^2}{2}.
\]
所以 $v^2-v_0^2=2\int_{x_0}^{x}a(\xi)\dd\xi$。这并未假设加速度恒定，而是把运动分支上的时间参数换成了位置参数。

积分结果给出的是速度平方，仍需检查根号内是否非负，以及所求位置能否在给定分支上到达。若速度在端点为零，位置参数在该处可能失去单调性，转向后的运动应另行分析。使用该式的前提是沿直线的带符号加速度与坐标方向一致；不能将加速度大小未经判断就当作正数代入。

\section{过河、追及与最小值}
\subsection{参考系选择怎样简化未知量}
选择参考系的目的，是使某些位置、速度或边界变简单。观察两车追及时，随前车平动可把相对位移直接与初始间距比较；分析车流间隙时，随车流匀速运动可把移动的间隙变成固定通道。换系之后仍需保留各物体相对该系的初始速度，不能只把目标移到静止而忘记观察者本身也改变了速度。

对两个共同受重力加速度的抛体，转到随共同自由落体运动的平动坐标中，它们的相对加速度为零，因此相对轨迹是直线。这里可用共同加速度相减求相对运动，但不能因此把该加速坐标系称为惯性系；本讲只用运动学变换，不在其中直接套无附加项的动力学定律。

\subsection{最短时间与最短路程}
设水速沿 $x$为 $u$，静水船速为 $w$，河宽为 $d$，船头偏离河岸法线向上游的角为 $\beta$。船对岸分量是 $V_x=u-w\sin\beta$、$V_y=w\cos\beta>0$，故 $T=d/V_y$，位移长度为 $S=d\sqrt{V_x^2+V_y^2}/V_y$。这样同一组分量已把两种优化目标明确分开。

图\ref{k2t-river-coach}为 $w<u$的速度圆，原点 $O$在圆外。要使 $S$最小，就要使合速度与河岸的夹角最大，因此 $OV$与速度圆相切。直角三角形 $OWV$给 $|OV|=\sqrt{u^2-w^2}$、$V_y=w\sqrt{u^2-w^2}/u$，于是 $S_{\min}=du/w$，对应 $\sin\beta=w/u$。这里并没有把水流抵消；船仍顺流漂移，只是漂移与过河速度的比值最小。
\fig{k2t-river-coach}{水速大于船速时，最短航程对应速度圆的切线}{\begin{tikzpicture}
\coordinate(O)at(0,0);\coordinate(W)at(3,0);\coordinate(V)at(2.52,1.1);
\draw[axis](0,0)--(4.7,0)node[right]{$V_x$};
\draw[axis](0,0)--(0,2)node[above]{$V_y$};
\draw[world](W)circle(1.2);\draw[flow](O)--(W)node[midway,below]{$u$};
\draw[flow,blue](W)--(V)node[midway,above right]{$w$};
\draw[flow,red](O)--(V)node[midway,above left]{$\bm V$};
\fill(O)circle(1.2pt)node[below left]{$O$};
\fill(W)circle(1.2pt)node[below right]{$W$};
\fill(V)circle(1.2pt)node[above left]{$V$};
\end{tikzpicture}}

船对水速率固定时，船对岸的速度端点位于以水流速度端点为圆心、船速为半径的圆上。最短过河时间要求垂直河岸的速度分量最大，因此船头应垂直河岸；最短对岸路程却要求合速度尽可能垂直河岸，二者目标不同。

若船速大于水速，能用向上游的分量抵消水流，走垂直河岸的直线。若船速小于水速，不能抵消水流，应从原点作速度圆的切线，以取得合速度与河岸的最大夹角。船速恰等于水速时，垂直合速度的极限同时意味着过河速度趋零，最短路程仅有下确界，不能说有限时间内真正达到该极限。这些分支应在求角度前先判断。

\subsection{追及应检查相对位置的最小值}
把后物体相对前物体的净间距记为 $D(t)$，则 $\dot D=v_{\text{前}}-v_{\text{后}}$。若存在 $\dot D=0$ 的时刻，它只是距离的驻点候选；还要核对是否属于当前运动阶段，并与阶段端点、停车时刻等比较。车辆在不同时间停下时，求得一个“速度相等时刻”后直接结束分析，可能漏掉停车后的继续接近。

避免相撞要求 $D(t)\ge0$ 对整个有效过程成立；若要求严格不接触，应使用严格不等式。临界条件常由 $D=0$ 与 $\dot D=0$ 同时给出，但当最危险状态位于终点时，只需端点条件。碰撞会突然改变速度，碰撞后的过程必须重新取初值；不能把碰前的相对运动抛物线继续延伸到碰后。
\subsection{最近距离与全程安全的统一判断}
若两物体相对位矢为 $\bm R(t)=\bm R_0+\bm U t$，最近距离不必由两条轨迹的交点判断。先作点积，$\bm R_0,\bm U$均为常矢量；在 $U\ne0$时再配方：
\begin{align*}
 R^2&=(\bm R_0+\bm U t)\cdot(\bm R_0+\bm U t)
 =R_0^2+2(\bm R_0\cdot\bm U)t+U^2t^2\\
 &=U^2\left[t^2+2\frac{\bm R_0\cdot\bm U}{U^2}t
 +\frac{(\bm R_0\cdot\bm U)^2}{U^4}\right]
 +R_0^2-\frac{(\bm R_0\cdot\bm U)^2}{U^2}\\
 &=U^2\left(t+\frac{\bm R_0\cdot\bm U}{U^2}\right)^2
 +R_0^2-\frac{(\bm R_0\cdot\bm U)^2}{U^2}.
\end{align*}
也可对展开式逐项求导：$\dd(R^2)/\dd t=2\bm R_0\cdot\bm U+2U^2t$，令其为零，得到同一顶点时刻。
当 $U\ne0$且驻点 $t_*=-\bm R_0\cdot\bm U/U^2$在考察区间内，它给最近距离。其几何意义是该瞬间连线垂直相对速度；若 $t_*<0$，未来最小距离出现在初始时刻。若只考察有限时间，还要比较终点。$U=0$时距离恒定，应单独处理，不能除以零。

制动题中相对加速度通常非零，但思想完全相同。令两车有向间距为 $D=x_{\text{前}}-x_{\text{后}}$，则 $\dot D=v_{\text{前}}-v_{\text{后}}$。每个制动阶段内，最小间距只可能出现在该阶段端点或速度相等处；反应结束、前车停车、后车停车会使阶段切换，也必须纳入比较。最终停车位置只是这些候选状态之一。

图\ref{k2t-gap-coach}给出一次内部最危险的示意：起初后车较快使 $D$下降，到 $t_*$两车同速，随后间距回升。以临界情况 $D(t_*)=0$求参数后，还应检查 $t_*$是否真的位于两车均在制动的区间。如果算出的同速时刻晚于前车停车，就应换用前车静止的下一段方程。正是这个区间检查产生了制动答案的不同分支。
\fig{k2t-gap-coach}{间距的极小值由全程运动而非仅由停车位置确定}{\begin{tikzpicture}
\draw[axis](0,0)--(5.8,0)node[right]{$t$};\draw[axis](0,0)--(0,3)node[above]{$D$};
\draw[world]plot[domain=0:5,samples=50]({\x},{.65+.27*(\x-2.5)^2});
\draw[guide](2.5,0)--(2.5,.65);\draw[blue](1.7,.65)--(3.3,.65);
\node[below]at(2.5,0){$t_*$};\node[above]at(2.5,.8){两车同速};
\end{tikzpicture}}

若题目仅考察 $0\le t\le T$，应先把 $t_*$与区间比较：$t_*<0$时距离从开始就增大，最小值在零时刻；$t_*>T$时考察结束前仍在接近，最小值在 $T$；只有 $t_*\in[0,T]$时，才能采用垂足位置。给定两物体的共同加速度时，相对加速度为零，同一判断仍成立。




图\ref{k2deep:closest}在相对位置空间中画出可考察的线段 $P_0P_T$，虚线则是把相对运动向过去和未来延长的结果。由原点 $O$ 作垂线得到 $H$，只有 $H$ 落在实线段上才是本题的最近位置；否则应取离 $H$ 较近的端点。这里最短距离是点到相对轨迹段的距离，并不是两条地面轨迹之间任取两点的距离。
\fig{k2deep:closest}{最近点必须属于实际考察的相对轨迹段}{\begin{tikzpicture}[scale=.85]
\foreach \shift/\lo/\hi/\title in {0/-1.2/1.2/{驻点在区间内},4/0.6/2/{驻点在过去},8/-2/-.6/{驻点在终点以后}}{
\begin{scope}[xshift=\shift cm]
\draw[guide](-2.2,1.5)--(2.2,1.5);\draw[blue,very thick](\lo,1.5)--(\hi,1.5);\draw[guide](0,0)--(0,1.5);
\fill(0,0)circle(1.5pt)node[below]{$O$};\draw[fill=white](0,1.5)circle(1.5pt)node[above]{$H$};
\fill(\lo,1.5)circle(1.5pt)node[below]{$P_0$};\fill(\hi,1.5)circle(1.5pt)node[below]{$P_T$};
\node at(0,-.65){\title};\end{scope}}
\end{tikzpicture}}


\subsection{最优值与临界接触}
临界安全距离常允许数学上“刚好接触而不穿越”。若题目说避免接触，就应在临界值之外留严格余量；若只说不相撞而把质点接触视作临界界限，应在答案中说明这一约定。最小值与下确界也不同：存在可行参数取得的才是最小值，只能无限接近但取不到的应称下确界。

优化某一辆车的制动时，其他车辆的策略也是条件的一部分。恒减速度、延迟开始制动和任意时变制动具有不同的可行集合。先明确策略，再证明候选策略对后车最有利，最后检查全程间距，才能完成优化证明，不能只比较最终停车位置就跳过途中最危险的状态。

% thirty-k2-theory
